% Section 18.13. The lesson's bilaga A has no summary of its own: the table collects the formulas
% of sections 18.3 to 18.10 with their C, and the rest is from lectures/L18/README.md: the
% objectives, the evaluation questions and the next lesson.

\section{Sammanfattning}
\label{c18:sec:sammanfattning}

\begin{booktable}{@{}>{\raggedright\arraybackslash}p{8.5em}p{12em}L@{}}
  \thead{Begrepp} & \thead{Matematik} & \thead{I C} \\
  \midrule
  Rot, potens & $\sqrt{x}$, $x^y$ & \code{sqrt(x)}, \code{pow(x, y)} \\
  Logaritmer & $\ln x$, $\lg x$ & \code{log(x)}, \code{log10(x)} \\
  Decibel & $20\log_{10}(U_{\text{ut}}/U_{\text{in}})$ & \code{20.0 * log10(u_out / u_in)} \\
  Grader till radianer & $\text{grader} \cdot \pi/180$ & \code{deg * PI / 180.0} \\
  Vinkel i rätt kvadrant & $\arg(x + jy)$ & \code{atan2(y, x)} \\
  Centraldifferens & $\bigl(f(x + h) - f(x - h)\bigr)/2h$ & \code{derivative(f, x, 1e-6)} \\
  Trapetsmetoden & $\int_a^b f(x)\,dx$, $n$ delintervall & \code{integral(f, a, b, n)} \\
  Komplext tal & $3 + 4j$ & \code{3.0 + 4.0 * I} \\
  Belopp, argument & $\abs{z}$, $\arg(z)$ & \code{cabs(z)}, \code{carg(z)} \\
  Polär form & $r e^{j\theta}$ & \code{r * cexp(I * theta)} \\
  Sampel & $\abs{U}\sin(2\pi k/N + \delta)$, $N = f_s/f$ &
    \code{amplitude * sin(2.0 * PI * k / n + delta)} \\
\end{booktable}

\needspace{6\baselineskip}
\subsection*{Det här ska du kunna}
Efter det här kapitlet ska du:
\begin{itemize}
  \item Förstå skillnaden mellan heltal och flyttal i C, och kunna välja rätt typ och
    formatspecificerare.
  \item Kunna använda \file{math.h} för att beräkna potenser, rötter, logaritmer och decibel.
  \item Kunna räkna med vinklar i radianer och omvandla till och från grader.
  \item Kunna beräkna en derivata numeriskt och jämföra med den analytiska derivatan.
  \item Kunna beräkna en bestämd integral numeriskt med trapetsmetoden.
  \item Kunna representera komplexa tal och fasorer med \file{complex.h} samt beräkna belopp och
    argument.
  \item Kunna beräkna sampelvärden för en sinussignal och skriva ut dem för vidare analys.
\end{itemize}

\testadigsjalv
\begin{enumerate}
  \item Vad blir resultatet av \code{1 / 2} i C, och varför?
  \item Vad gör flaggan \code{-lm}, och vilket felmeddelande får man om matematikbiblioteket inte
    länkas?
  \item Varför ger \code{sin(30)} inte $0{,}5$?
  \item Hur beräknas derivatan av $f(x)$ i punkten $x$ numeriskt, och vad styr noggrannheten?
  \item Hur beräknas beloppet och argumentet för ett komplext tal i C?
\end{enumerate}

\needspace{6\baselineskip}
\subsection*{Nästa kapitel}
\Kapref{19} avslutar boken med en repetition av kursens innehåll och förberedelser inför
tentamen.
